\documentclass[aspectratio=169,11pt]{beamer}

\usetheme{Madrid}
\usecolortheme{default}
\usefonttheme{professionalfonts}
\setbeamertemplate{navigation symbols}{}
\setbeamertemplate{items}[circle]

\usepackage[T1]{fontenc}
\usepackage[utf8]{inputenc}
\usepackage[polish]{babel}
\usepackage{lmodern}
\usepackage{amsmath}
\usepackage{booktabs}
\usepackage{array}
\usepackage{tikz}
\usepackage{pgfplots}
\pgfplotsset{compat=1.18}
\usetikzlibrary{arrows.meta, positioning, shapes.geometric, fit, backgrounds}

\definecolor{MainBlue}{RGB}{27,78,114}
\definecolor{Accent}{RGB}{204,102,0}
\definecolor{Soft}{RGB}{241,245,249}
\definecolor{Good}{RGB}{46,125,50}
\definecolor{Warn}{RGB}{183,28,28}

\setbeamercolor{structure}{fg=MainBlue}
\setbeamercolor{frametitle}{fg=white,bg=MainBlue}
% \setbeamercolor{title}{fg=MainBlue}
\setbeamercolor{block title}{fg=white,bg=MainBlue}
\setbeamercolor{block body}{bg=Soft}
\setbeamercolor{alerted text}{fg=Accent}
% \setbeamertemplate{title page}{%
%   \vspace*{0.8cm}
%   {\usebeamerfont{title}\color{MainBlue}\inserttitle\par}
%   \vspace{0.25cm}
%   {\large \insertsubtitle\par}
%   \vspace{0.8cm}
%   {\normalsize \insertauthor\par}
%   \vspace{0.15cm}
%   {\small \insertdate\par}
%   \vfill
%   \begin{beamercolorbox}[wd=\textwidth,sep=0.5em,left]{block body}
%   \textbf{Cel prezentacji:} syntetyczne omówienie analizy drzewa zdarzeń z rozmytymi prawdopodobieństwami dla systemu ochrony przeciwpożarowej magazynu logistycznego.
%   \end{beamercolorbox}
% }

\title{Inteligencja obliczeniowa i jej zastosowania - Lista 2 - wariant B}
\subtitle{Drzewa zdarzeń z rozmytymi prawdopodobieństwami}
\author[Jakub Jaśków, Justyna Ziemichód]{Jakub Jaśków, Justyna Ziemichód}
\date{14 marca 2026}

\newcommand{\trap}[4]{$(#1,#2,#3,#4)$}
\newcommand{\smallcaps}[1]{{\footnotesize #1}}

\begin{document}

\begin{frame}
  \titlepage
\end{frame}

\begin{frame}{Plan prezentacji}
    \tableofcontents
\end{frame}

\section{Opis systemu}
\begin{frame}{Analizowany system i logika obrony}
\begin{columns}[T,totalwidth=\textwidth]
\column{0.57\textwidth}
\begin{block}{System}
Analizowanym obiektem jest \textbf{wielkopowierzchniowy magazyn logistyczny}, a zdarzeniem inicjującym jest \textbf{wybuch pożaru} o rozmytej intensywności $I_{IE}$ [1/rok].
\end{block}

\vspace{0.2cm}
\begin{itemize}
  \item \textbf{B1 -- zraszacze:} automatyczne tłumienie pożaru w zarodku.
  \item \textbf{B2 -- grodzie i drzwi ppoż.:} ograniczenie rozprzestrzeniania ognia.
  \item \textbf{B3 -- Straż Pożarna:} skuteczna interwencja w czasie krótszym niż 15 minut.
\end{itemize}

\column{0.4\textwidth}
\centering
\begin{tikzpicture}[node distance=0.4cm]
  \node[draw, rounded corners, fill=Soft, minimum width=2.8cm, minimum height=0.75cm, align=center] (ie) {Pożar\\$I_{IE}$};
  \node[draw, rounded corners, fill=blue!8, below=of ie, minimum width=2.8cm, minimum height=0.75cm, align=center] (b1) {B1\\Zraszacze};
  \node[draw, rounded corners, fill=blue!8, below=of b1, minimum width=2.8cm, minimum height=0.75cm, align=center] (b2) {B2\\Drzwi ppoż.};
  \node[draw, rounded corners, fill=blue!8, below=of b2, minimum width=2.8cm, minimum height=0.75cm, align=center] (b3) {B3\\Straż Pożarna};
  \node[draw, rounded corners, fill=orange!12, below=of b3, minimum width=2.8cm, minimum height=0.75cm, align=center] (c) {Konsekwencje\\C1--C4};
  \draw[-{Latex[length=2mm]}, thick, MainBlue] (ie) -- (b1);
  \draw[-{Latex[length=2mm]}, thick, MainBlue] (b1) -- (b2);
  \draw[-{Latex[length=2mm]}, thick, MainBlue] (b2) -- (b3);
  \draw[-{Latex[length=2mm]}, thick, MainBlue] (b3) -- (c);
\end{tikzpicture}

\vspace{0.35cm}
{\small \textcolor{Accent}{Model zakłada trzy niezależne warstwy zabezpieczeń.}}
\end{columns}
\end{frame}

\begin{frame}{Rozmyte skale wejściowe}
\begin{columns}[T,totalwidth=\textwidth]
\column{0.48\textwidth}
\begin{block}{Intensywność pożaru $I$ [1/rok]}
\centering
\small
\begin{tabular}{>{\raggedright\arraybackslash}p{1.5cm}p{4.0cm}}
\toprule
\textbf{Skala} & \textbf{Trapez} \\
\midrule
Bardzo niska & \trap{0.0}{0.0}{0.01}{0.02} \\
Niska & \trap{0.01}{0.02}{0.03}{0.05} \\
Średnia & \trap{0.03}{0.05}{0.08}{0.10} \\
Wysoka & \trap{0.08}{0.10}{0.15}{0.20} \\
\bottomrule
\end{tabular}
\end{block}

\column{0.48\textwidth}
\begin{block}{Prawdopodobieństwo awarii $P$}
\centering
\small
\begin{tabular}{>{\raggedright\arraybackslash}p{1.5cm}p{4.0cm}}
\toprule
\textbf{Skala} & \textbf{Trapez} \\
\midrule
Bardzo niskie & \trap{0.0}{0.0}{0.001}{0.005} \\
Niskie & \trap{0.001}{0.005}{0.010}{0.050} \\
Średnie & \trap{0.010}{0.050}{0.100}{0.150} \\
Wysokie & \trap{0.150}{0.300}{1.0}{1.0} \\
\bottomrule
\end{tabular}
\end{block}
\end{columns}

\end{frame}

\section{Drzewo zdarzeń}
\begin{frame}{Drzewo zdarzeń}
\begin{figure}
    \centering
    \includegraphics[width=0.7\linewidth]{drzewo\_zdarzeń.png}
    \label{fig:placeholder}
\end{figure}
\end{frame}

\section{Oszacowanie parametrów modelu}
\begin{frame}{Oszacowanie parametrów modelu}
\small
\begin{columns}[T,totalwidth=\textwidth]
\column{0.58\textwidth}
\begin{block}{Przyjęte wartości rozmyte}
\begin{tabular}{>{\raggedright\arraybackslash}p{3.2cm}p{3.6cm}}
\toprule
\textbf{Parametr} & \textbf{Wartość} \\
\midrule
Intensywność pożaru $I_{IE}$ & Średnia = \trap{0.03}{0.05}{0.08}{0.10} \\
Awaria B1: $P_{B1}$ & Niska = \trap{0.001}{0.005}{0.010}{0.050} \\
Awaria B2: $P_{B2}$ & Bardzo niska = \trap{0.0}{0.0}{0.001}{0.005} \\
Awaria B3: $P_{B3}$ & Średnia = \trap{0.010}{0.050}{0.100}{0.150} \\
\bottomrule
\end{tabular}
\end{block}

\column{0.4\textwidth}
\begin{block}{Sukces barier}
\[
1 \ominus P_{awarii}=(1-d,1-c,1-b,1-a)
\]
\vspace{-0.25cm}
\begin{align*}
1\ominus P_{B1} &= (0.950,0.990,0.995,0.999) \\
1\ominus P_{B2} &= (0.995,0.999,1.0,1.0) \\
1\ominus P_{B3} &= (0.850,0.900,0.950,0.990)
\end{align*}
\end{block}
\end{columns}

\begin{alertblock}{Założenie obliczeniowe}
Dla intensywności konsekwencji zastosowano aproksymację mnożenia trapezów:
\[
A\otimes B=(a_A a_B,\; b_A b_B,\; c_A c_B,\; d_A d_B)
\]
\end{alertblock}
\end{frame}

\begin{frame}{Wyniki intensywności konsekwencji C1--C4}
\small
\begin{table}
\centering
\scriptsize
\begin{tabular}{>{\raggedright\arraybackslash}p{1.3cm}>{\raggedright\arraybackslash}p{5.3cm}p{4.7cm}}
\toprule
\textbf{Konsekw.} & \textbf{Ścieżka} & \textbf{$I_C$ [1/rok]} \\
\midrule
C1 & $I_{IE}\otimes (1\ominus P_{B1})$ & \trap{0.0285}{0.0495}{0.0796}{0.0999} \\
C2 & $I_{IE}\otimes P_{B1}\otimes (1\ominus P_{B2})$ & \trap{0.000029}{0.000249}{0.000800}{0.005000} \\
C3 & $I_{IE}\otimes P_{B1}\otimes P_{B2}\otimes (1\ominus P_{B3})$ & \trap{0.0}{0.0}{0.00000076}{0.00002475} \\
C4 & $I_{IE}\otimes P_{B1}\otimes P_{B2}\otimes P_{B3}$ & \trap{0.0}{0.0}{0.00000008}{0.00000375} \\
\bottomrule
\end{tabular}
\end{table}

\vspace{0.3cm}

\begin{alertblock}{Najważniejsza obserwacja}
Dominującym scenariuszem jest \textbf{C1} — najczęściej pożar zostaje opanowany na pierwszej warstwie ochrony.  
Konsekwencja \textbf{C4} ma intensywność rzędu $10^{-8}$–$10^{-6}$ zdarzeń na rok.
\end{alertblock}

\end{frame}

\begin{frame}{Funkcje trapezowe intensywności konsekwencji}

\begin{center}
\begin{tikzpicture}

\begin{axis}[
width=12cm,
height=8cm,
xmode=log,
log basis x=10,
xmin=1e-8,
xmax=2e-1,
ymin=0,
ymax=1.1,
xlabel={Intensywność zdarzenia $I_C$ [1/rok]},
ylabel={Stopień przynależności},
grid=major,
minor grid style={gray!15},
grid style={gray!40},
legend pos=outer north east,
clip=false,
tick label style={font=\scriptsize},
label style={font=\small}
]

% C1
\addplot[thick,blue]
coordinates{
(0.0285,0)
(0.0495,1)
(0.0796,1)
(0.0999,0)
};
\addlegendentry{C1}

% C2
\addplot[thick,red]
coordinates{
(0.000029,0)
(0.000249,1)
(0.000800,1)
(0.005000,0)
};
\addlegendentry{C2}

% C3
\addplot[thick,green!60!black]
coordinates{
(1e-8,0)
(1e-8,1)
(0.00000076,1)
(0.00002475,0)
};
\addlegendentry{C3}

% C4
\addplot[thick,orange]
coordinates{
(1e-8,0)
(1e-8,1)
(0.00000008,1)
(0.00000375,0)
};
\addlegendentry{C4}

\end{axis}

\end{tikzpicture}
\end{center}

\end{frame}

\section{Analiza znaczenia i wrażliwości}
\begin{frame}{Analiza znaczenia i wrażliwości dla katastrofy C4}
\small
\begin{columns}[T,totalwidth=\textwidth]
\column{0.48\textwidth}
\begin{block}{Analiza znaczenia}
\[
I_{C4}=I_{IE}\otimes P_{B1}\otimes P_{B2}\otimes P_{B3}
\]
\begin{itemize}
  \item Ścieżka do katastrofy ma \textbf{charakter szeregowy}.
  \item Wyeliminowanie awarii \emph{dowolnej} bariery daje:
  \[
  I_{C4}= (0,0,0,0)
  \]
  \item Oznacza to, że idealna niezawodność jednej warstwy całkowicie blokuje scenariusz C4.
\end{itemize}
\end{block}

\column{0.5\textwidth}
\begin{block}{Analiza wrażliwości}
\centering
\scriptsize
\begin{tabular}{>{\raggedright\arraybackslash}p{1.7cm}>{\centering\arraybackslash}p{1.2cm}>{\raggedright\arraybackslash}p{3.2cm}}
\toprule
\textbf{Zmiana} & \textbf{Skala} & \textbf{Nowe $I_{C4}$} \\
\midrule
Stan bazowy & -- & $\left(0,0,8.0\cdot10^{-8},\;3.75\cdot10^{-6}\right)$ \\
B1 gorsze & N $\to$ S & $\left(0,0,8.0\cdot10^{-7},\;1.12\cdot10^{-5}\right)$ \\
B3 lepsze & S $\to$ N & $\left(0,0,8.0\cdot10^{-9},\;1.25\cdot10^{-6}\right)$ \\
\bottomrule
\end{tabular}
\end{block}

\vspace{0.1cm}
\begin{alertblock}{Wniosek}
Najsilniej warto dbać o \textbf{wczesne warstwy ochrony}, zwłaszcza zraszacze -- ich pogorszenie podnosi intensywność katastrofy około dziesięciokrotnie.
\end{alertblock}
\end{columns}
\end{frame}

\section{Wnioski}

\begin{frame}{Wnioski końcowe}

\begin{block}{Wnioski z analizy}
\begin{enumerate}
\item Model o strukturze iloczynowej jest wrażliwy na zmiany parametrów niezawodności barier bezpieczeństwa.

\item Pogorszenie niezawodności zraszaczy z poziomu \textbf{Niskiego} do \textbf{Średniego} powoduje około \textbf{10-krotny wzrost} intensywności scenariusza katastrofalnego  
($8.0 \times 10^{-8} \rightarrow 8.0 \times 10^{-7}$).

\item Największe znaczenie dla ograniczenia ryzyka mają \textbf{pierwsze warstwy ochrony}.
\end{enumerate}
\end{block}

\vspace{0.2cm}

\begin{alertblock}{Implikacja praktyczna}
Kluczowe dla bezpieczeństwa obiektu jest utrzymanie wysokiej niezawodności systemu zraszaczy.
\end{alertblock}

\end{frame}

\begin{frame}[plain]

\begin{center}

\vspace{1.5cm}

{\Huge \textcolor{MainBlue}{Dziękujemy za uwagę}}


\end{center}

\end{frame}

\end{document}
